% PROPOSED supplement section: the technical details the six-figure wrap moved out of the Methods and CLAIM 2024
% still asks for (items 9, 13, 22-25, 27; packet_src/claim2024/map.json). NOT \input anywhere yet: the edit phase,
% which owns supplement.tex, decides whether to take it (% PROPOSED supplement section: the technical details the six-figure wrap moved out of the Methods and CLAIM 2024
% still asks for (items 9, 13, 22-25, 27; packet_src/claim2024/map.json). NOT \input anywhere yet: the edit phase,
% which owns supplement.tex, decides whether to take it (% PROPOSED supplement section: the technical details the six-figure wrap moved out of the Methods and CLAIM 2024
% still asks for (items 9, 13, 22-25, 27; packet_src/claim2024/map.json). NOT \input anywhere yet: the edit phase,
% which owns supplement.tex, decides whether to take it (% PROPOSED supplement section: the technical details the six-figure wrap moved out of the Methods and CLAIM 2024
% still asks for (items 9, 13, 22-25, 27; packet_src/claim2024/map.json). NOT \input anywhere yet: the edit phase,
% which owns supplement.tex, decides whether to take it (\input{packet_src/supplement_technical} after
% \input{tables/tab2_detection}, so its table numbers S3). The supplement has no word limit (12 pages).
% Every value carries its source in the comment beside it;
% \PENDING marks what no file on disk yet states. Code SHA is absent from the level-model result files
% (docs/MODEL_CARD.md 1.2), so the training settings are those of the scripts named, at the revision on disk on
% 2026-09-23; paper 2's manifest restates them for the split-trained arm (the comparison is paper 2 since 2026-09-24,
% PAPER1_SCOPE.md).

\section*{Technical details}

\paragraph{Pre-processing.}
% docs/REAL_XRAY_DATASETS_SPEC.md "Display" and "Boxes"
Every real exposure was converted to attenuation (the negative logarithm of a linear detector signal against its
air or flood value, or a linear reading of a stored logarithmic display), then windowed identically for every
device: the field's 1st, 5th, 50th, 95th and 99th percentiles were mapped to grey levels 27, 47, 126, 216 and 239,
bone bright, and stored as an 8-bit image. A truth box is the projection of every voxel of a cone-beam CT vertebra
label (subsampled to at most 20,000 points) through the registered pose, taking the 2nd--98th percentile extent;
a box was kept when at least half its projected points fell inside the image and both of its sides measured at
least 8 pixels. Level crops are fixed-scale two-channel windows of 320 pixels (image and target marker) with the
target body at 96 pixels.

\paragraph{Acquisition.}
% docs/RSNAI_TABLE1_DATASETS.md (pitch, frame size, stored value, tube voltage p5-p95); docs/MODEL_CARD.md 6.1
% (C-ArmLSpine source-image and source-object distance). Retrospective public releases: no acquisition protocol
% beyond what the headers carry; "not recorded" where the release gives none.
Table~\ref{tab:acquisition} lists the acquisition parameters each release carries. No release documents an exposure protocol beyond
its headers, and no parameter was controlled by this study.

\begin{table}[tbh!]
\centering
\caption{Acquisition parameters of the real exposures, as released.\label{tab:acquisition}}
\footnotesize
\begin{tabular}{>{\raggedright\arraybackslash}p{2.5cm}>{\raggedright\arraybackslash}p{3.6cm}>{\raggedright\arraybackslash}p{1.7cm}>{\raggedright\arraybackslash}p{3.3cm}>{\raggedright\arraybackslash}p{2.6cm}}
\toprule
Collection & Device & Detector pitch (mm) & Frame (pixels) & Tube voltage, kVp (5th--95th percentile) \\
\midrule
C-ArmLSpine & Surgio B51 flat-panel C-arm & 0.15 (a) & 2,048 $\times$ 2,048; exported at 1,024 & 73--111 \\
DeepFluoro & Siemens Cios Fusion & 0.194 & 1,536 $\times$ 1,536 & not recorded (b) \\
COBRA2026 & Elekta XVI kV imager & 0.8 & 504 $\times$ 504 & 120 \\
Geneva hip & Philips BV Pulsera image intensifier & 0.2904 & 1,000 $\times$ 1,000 & 65--90 \\
Ljubljana 2D-3D-GS & Trixell PIXIUM 4600 flat panel & 0.286 & 1,560 $\times$ 1,560 & not recorded \\
\bottomrule
\end{tabular}

\smallskip
\raggedright\footnotesize
(a) 307.2 mm over 2,048 pixels; the header pixel spacing is a placeholder of 1.0. Source-to-image distance
1,140 mm, source-to-object distance 745 mm, object-plane pixel 0.1915 mm at the 1,024-pixel export.
(b) Present in the vendor DICOM headers of the same acquisitions but not joined per registered frame.
\end{table}

\paragraph{Models, initialisation and training.}
% Held-out-device detector: atlas/model.json "Architecture" (args.yaml and data.yaml of rx_lodo_carmlspine_1024 read
% on the grid 2026-09-22); results/a1.json operating_points/cross_device_0.74 (how_chosen).
The sacrum detector is YOLO11m (Ultralytics; classes vertebra and sacrum; 1,024-pixel input), initialised from the
COCO-pretrained YOLO11m weights through a render-trained checkpoint and trained for 80 epochs (patience 25, batch
8, seed 0, deterministic) with rotation up to 10\textdegree, translation 0.1, scale 0.3, horizontal flip 0.5 and
mosaic 0.5, on the training and validation frames of COBRA2026 and DeepFluoro; its checkpoint was selected on
their test frames, so only its held-out-device results are clean.
% Keypoint detector: docs/MANIFEST_DATA_MODELS_RESULTS.md row detector_v5z_final; Methods (confidence 0.35);
% chain detect_v5z_final <- detect_v3 <- public ultralytics/yolo11m-pose.pt (COCO-pose), results/eval_gate/LINEAGE.md.
The body detector is YOLO11m-pose fine-tuned for 60 epochs on the plain-film keypoint corpus from an earlier pose
checkpoint (detect\_v3), itself trained from the COCO-pretrained YOLO11m-pose weights, and used at confidence 0.35
\PENDING{G5}{keypoint detector: training corpus by source and split, and whether any film of a test person is in it}.
% Encoder: docs/MODEL_CARD.md 1.1 and 6.3 (architecture, 11,173,376 parameters, 3.62 GMAC); slurm/carmlspine_cv.sh
% lines 21 and 31 (init alias allspine_v2_r1, 15 epochs, lr 1e-4, --augment 2); tools/train_level_classifier.py
% (AdamW, weight decay 1e-4, cosine schedule, class-weighted cross-entropy).
The level encoder is ResNet-18 with a two-channel stem and a 27-level head (11,173,376 parameters; 3.62 giga
multiply--accumulates per crop). The cross-validation encoders were initialised from a level encoder trained on
renders and films (registry model allspine\_v2\_r1) and fine-tuned per fold for 15 epochs at a learning rate of
$10^{-4}$ with AdamW (weight decay $10^{-4}$), a cosine schedule and class-weighted cross-entropy.
% Figure 2 encoder: results/figure2/figure2_metrics.json "encoder", "encoder_sha256"; ImageNet root per
% results/eval_gate/LINEAGE.md (allspine_v9 <- public torchvision resnet18 IMAGENET1K_V1, "used when cached on the
% node; random init otherwise"). Checked 2026-09-24 (grid jobs 40429822, 40429828, read-only): the training script
% loads the ImageNet file when it is in the torch hub cache and prints "pretrained=<bool>"; the file was in the shared
% home cache from 2026-09-21 18:31, and v9 trained from 2026-09-22 21:57 (model_night/logs/v9train_40415669.out), but
% that log starts at epoch 25, so the load itself is not recorded. The sentence says so.
% Checkpoint selection: tools/train_level_classifier.py keeps the best epoch by exact accuracy on the bank's validation
% split (results/eval_gate/EVALUATORS.md item 1), which is why Table 3 (tab3_naming.tex; Table 4 before 2026-09-24
% evening) labels the film rows "selection split".
The encoder whose embedding Figure~2 maps is the allspine\_v9 model (SHA-256 188f9fcb\ldots), started by the training
script's default from ImageNet weights, which were cached on the computing grid before it trained (its log does not
record the load); like every level encoder here, its checkpoint was the epoch with the best exact accuracy on its
bank's validation split, so numbers on that split are selection-split numbers.
% Head: docs/MODEL_CARD.md 1.1 (RunLabeller, 2,117,403 parameters); tools/head_mixed_length.py (defaults epochs 30,
% lr 1e-3, AdamW weight decay 1e-4, cosine; per-position cross-entropy); slurm/carmlspine_cv.sh line 33
% (--train-mode mixed, no epoch or learning-rate override).
The recurrent head (linear 512 to 256, GELU, two-layer bidirectional GRU of width 256, linear 512 to 27; 2,117,403
parameters) was trained on runs of two to nine consecutive embeddings of mixed length for 30 epochs at a learning
rate of $10^{-3}$ (AdamW, weight decay $10^{-4}$, cosine schedule) with per-position cross-entropy.
% docs/MODEL_CARD.md 6.3 ("flip TTA on"); tools/alphabet_couple.py header (posterior "flip-TTA averaged"); no ensemble.
No ensemble is used; each crop is also read horizontally flipped and the two readings averaged.
Counting and the alphabet coupling have no learned parameters.
% Split-trained arm: docs/RESUME_2026-09-22_PM.md checkpoint 2026-09-23 21:45 (the v8/v9 banks trained on 49 of 54 L6
% validation and test persons of the built split, tables/tab4_transitional.md note d, so the arm starts from
% ImageNet on the split's training side). Since 2026-09-24 (evening) that arm is scored in paper 2 (the former
% "Figure 6 arm"; paper2_levelcheck/); paper 1 shows only its development learning curve (supplement section
% "Development learning curve of a split-trained level model"), so the sentence names no scored configuration.
A level model trained from ImageNet weights on the training side of the person-disjoint transitional split,
because the banks behind the encoders above hold most of its validation and test persons, is reported here only
through its development learning curve; no model in this manuscript is scored on a test person of that split.
 after
% % SUPPLEMENT Table S2 since the 2026-09-23 six-figure wrap (formerly main Table 2); input by supplement.tex.
% Table 2 -- detection per held-out device and training arm at both operating points, with case-clustered intervals
% and PR-AUC. No merged cells: the target block heading sits in the first cell of its first row. Sources, one line per
% cell, in tab2_detection.sources.md; the same rows with BCa intervals, plus the arms cut from the page for space
% (the design-A seed pairs, renders-only on COBRA and DeepFluoro, Cios-only on COBRA, Elekta-only and Surgio-only on
% DeepFluoro), in tab2_detection.md.
%   results/a1.json                  the reported detector (lodo_carm) on C-ArmLSpine at 0.74 and 0.25 (k/n, percentile, BCa)
%   results/lcurve/lcurve.json       design-B stage-1 arms and the same-device controls at 0.74 and 0.25, PR-AUC (both devices)
%   results/contribution_matrix.json every other arm at 0.25 (recall as a fraction; k = recall x n, exact integers)
%   results/a3.json                  the plain-film block (render-trained vs film-trained, conf 0.05 and 0.25)
%   results/fig3_detection/scored/*  the 0.74 cells of the render-only and real-only held-out arms (same sweeps, both points)
% BOX CONVENTION (A28): every number here is in the whole-vertebra convention, and the sentence in the text that quotes a
% row must say "whole-vertebra". The vertebral-body-box row per block (the all-device detector det_g3_primary_s0, 15 G5
% cells) was DROPPED from paper 1 on 2026-09-24 (Greg): it can be scored only on the held-out test persons paper 2 keeps
% blinded; the rows are kept verbatim in paper2_levelcheck/MOVED_FROM_PAPER1.md (item 10).
% Labels: tab:detection is the table; tab:contribution is kept as an alias so the Results prose that still cites the
% retired contribution table resolves to this one (the single-source arms it quoted are rows here).
% Revision 2 (2026-09-24, REVISION_LOG2.md): plain-radiograph rows re-scored on the 17 NHANES II lumbar spines neither
% detector trained or was selected on (EVAL_LEAKS L2; closure MANUSCRIPT_CORRECTIONS.md 3a;
% results/eval_gate/leaks/render_vs_film_unseen_rerun.json); zero-floor percentile cells given unclustered Wilson
% intervals (REVIEW2_STATS N18; results/stats_review2/recompute.json#S10_floor_cells).
% Layout: landscape (rotated 23.2 cm minipage, see tab_corpus.tex); body 7 pt on 7.6 pt.
\begin{table}[p]
\centering
\begingroup
\singlespacing
\fontsize{7}{7.6}\selectfont
\newcommand{\fnm}[1]{\textsuperscript{\fontsize{7}{7}\selectfont #1}}
\setlength{\tabcolsep}{2pt}
\renewcommand{\arraystretch}{1.0}
\rotatebox[origin=c]{90}{\begin{minipage}{23.2cm}
\captionsetup{font=footnotesize,skip=4pt}
\centering
\caption{Detection per held-out device and training arm at the cross-device (0.74) and default (0.25) operating points. Whole-vertebra box convention throughout; the all-device detector trained on vertebral-body boxes is reported in the separate held-out study.}
\label{tab:detection}\label{tab:contribution}
% widths sum to 22.00 cm; 14 gaps of 2 pt (0.99 cm) -> 22.99 cm
\begin{tabular}{%
>{\raggedright\arraybackslash}p{2.20cm}%
>{\raggedright\arraybackslash}p{5.20cm}%
>{\raggedright\arraybackslash}p{3.15cm}%
>{\raggedright\arraybackslash}p{3.00cm}%
>{\raggedright\arraybackslash}p{3.15cm}%
>{\raggedright\arraybackslash}p{3.00cm}%
>{\raggedright\arraybackslash}p{2.30cm}}
\toprule
Target device\fnm{a} & Training arm & Vertebra recall at 0.74\fnm{b} & Sacrum found at 0.74 & Vertebra recall at 0.25 & Sacrum found at 0.25 & PR-AUC\fnm{c}, vertebra / sacrum \\
\midrule
C-ArmLSpine (Surgio-B51) & Reported detector: real COBRA2026 and DeepFluoro frames, no render & 80.5 (231/287) [74.1, 86.8] & 76.8 (43/56) [64.3, 89.3] & 90.6 (260/287) [84.3, 94.5]\fnm{g} & 83.9 (47/56) [73.2, 92.9] & -- \\
 & Design B: styled renders and real frames of the two other devices; seed 1 & 63.1 (181/287) [54.8, 71.0] & 66.1 (37/56) [53.6, 78.6] & 87.8 (252/287) [81.4, 93.2] & 85.7 (48/56) [76.8, 92.9] & 0.890 / 0.900 \\
 & Elekta XVI only, styled and real & -- & -- & 86.8 (249/287) [81.1, 91.8] & 78.6 (44/56) [67.9, 89.3] & -- \\
 & Cios Fusion only, styled and real & -- & -- & 12.9 (37/287) [8.3, 18.1] & AP 28/28, lateral 0/28\fnm{g} & -- \\
 & Elekta XVI only, real frames cut to 240 in 10 patients & -- & -- & 39.0 (112/287) [33.2, 44.7] & 37.5 (21/56) [28.6, 46.4] & -- \\
 & Renders only, no real frame & 0.3 (1/287) [0.1, 1.9]\fnm{g} & 5.4 (3/56) [1.8, 14.6]\fnm{g} & 31.7 (91/287) [27.1, 36.1] & 53.6 (30/56) [44.6, 62.5] & -- \\
 & Control\fnm{d}: C-ArmLSpine real training frames only & 40.4 (116/287) [37.8, 42.9] & 94.6 (53/56) [85.4, 98.2]\fnm{g} & 41.5 (119/287) [38.8, 44.0] & 98.2 (55/56) [90.6, 99.7]\fnm{g} & 0.367 / 0.979 \\
 & Control\fnm{d}: real training frames of all three devices & 40.1 (115/287) [37.7, 42.3] & 96.4 (54/56) [87.9, 99.0]\fnm{g} & 40.8 (117/287) [38.4, 43.0] & 98.2 (55/56) [90.6, 99.7]\fnm{g} & 0.385 / 0.982 \\
\addlinespace[1.5pt]
COBRA2026 (Elekta XVI) & Real DeepFluoro and C-ArmLSpine frames, no render & 6.2 (55/893) [2.9, 9.4] & 3.4 (21/609) [1.3, 5.9] & 18.5 (165/893) [13.7, 22.8] & 13.3 (81/609) [8.1, 18.8] & -- \\
 & Design B: styled renders and real frames of the two other devices & 10.1 (90/893) [6.7, 13.1] & 13.0 (79/609) [9.0, 17.1] & 32.7 (292/893) [28.4, 36.2] & 36.5 (222/609) [31.6, 40.9] & 0.266 / 0.391 \\
 & Surgio-B51 only, styled and real; two runs of one list & -- & -- & 7.5 (67/893) [3.5, 12.1]; 12.2 (109/893) [8.0, 16.0] & 3.8 (23/609) [1.8, 6.2]; 8.0 (49/609) [5.3, 10.9] & -- \\
 & Control\fnm{d}: real training frames of all three devices & 93.2 (832/893) [86.4, 97.7] & 90.6 (552/609) [83.8, 96.3] & 97.4 (870/893) [93.3, 99.8] & 97.9 (596/609) [95.7, 99.5] & 0.963 / 0.985 \\
\addlinespace[1.5pt]
DeepFluoro (Cios Fusion)\fnm{e} & Real COBRA2026 and C-ArmLSpine frames, no render; test specimen & 74.7 (62/83)\fnm{e} & 100.0 (90/90)\fnm{e} & 95.2 (79/83)\fnm{e} & 100.0 (90/90)\fnm{e} & -- \\
 & Design B: styled renders and real frames of the two other devices; six specimens & -- & -- & 78.6 (338/430) [68.8, 92.0] & 95.0 (321/338) [78.5, 99.8] & -- \\
 & The same arm; test specimen & 69.9 (58/83)\fnm{e} & 87.8 (79/90)\fnm{e} & 94.0 (78/83)\fnm{e} & 100.0 (90/90)\fnm{e} & 0.923 / 0.954 \\
 & Control\fnm{d}: real training frames of all three devices; test specimen & 73.5 (61/83)\fnm{e} & 100.0 (90/90)\fnm{e} & 95.2 (79/83)\fnm{e} & 100.0 (90/90)\fnm{e} & 0.685 / 0.968 \\
\addlinespace[1.5pt]
Plain radiographs\fnm{f} & Render-trained detector, no real image & at 0.05: 67.7 (109/161) & -- & 34.8 (56/161) & -- & -- \\
 & Film-trained detector & at 0.05: 99.4 (160/161) & -- & 98.8 (159/161) & -- & -- \\
\bottomrule
\end{tabular}
\par\vspace{1pt}
\begingroup\fontsize{7}{7.6}\selectfont\raggedright
Cells are \% ($k/n$) [95\,\% CI]. CI = confidence interval; PR-AUC = area under the precision--recall curve; IoU = intersection over union. Greedy one-to-one matching at IoU $\ge$ 0.5; 95\,\% percentile bootstraps of 2,000 resamples clustered by case (C-ArmLSpine), patient (COBRA2026) or specimen (DeepFluoro), BCa intervals in the Markdown twin; ``--'' = not scored at that point or not computed. Styled = rendered through the fitted terms; design B = renders under each device's own poses.
\fnm{a}Test sets: C-ArmLSpine 56 frames, 28 cases, 287 vertebrae, 56 sacra; COBRA2026 648 frames, 27 patients, 893 vertebrae, 609 frames with a sacrum; DeepFluoro six specimens, 366 frames, 430 vertebrae, 338 with a sacrum; plain radiographs 17 NHANES II lumbar films (17 spines, 161 bodies) whose patients neither detector trained or was selected on. Truth boxes: C-ArmLSpine rows against whole-vertebra outline boxes projected from the CBCT labels, the convention every number on disk uses (under the SPINEPS vertebral-body boxes of the pending flagship rows the reported detector reads 5.2\,\% (15/287) at 0.25, a convention mismatch); COBRA2026 and DeepFluoro rows against whole-vertebra boxes projected from the registered CT. The device is held out unless the row says control.
\fnm{b}0.74 = best pooled vertebra F1 of the reported detector's sweeps on COBRA2026, DeepFluoro and the Geneva phantom, reading no C-ArmLSpine frame, adopted after the 0.25 result was reported; 0.25 = the software default; the reported detector chose its checkpoint on the COBRA2026 and DeepFluoro test splits.
\fnm{c}Non-interpolated average precision over the confidence sweep, computed for the design-B arms and the controls only.
\fnm{d}Controls trained on the target device's own frames (test cases held out by case); on C-ArmLSpine the outline convention penalises these body-trained arms, which read 89.9\,\% (258/287) and 90.6\,\% (260/287) at 0.25 under body boxes.
\fnm{e}One specimen is one cluster, no interval; the six-specimen rows test arms that never saw the device, their six-cluster intervals descriptive.
\fnm{f}Scored at 0.05 and 0.25 only, on every annotated NHANES~II lumbar spine that neither detector trained or was selected on; paired difference at 0.05 +31.7 percentage points [21.0, 45.4], sign test 16 vs 0 spines, Holm-adjusted $P<.001$; 17 clusters, so the intervals are approximate.
\fnm{g}Wilson 95\,\% interval where the percentile one reached 0 or 100\,\% against the data (unclustered for the two 0 cells) or the file flags errors in few clusters (sacrum controls unadjusted, design effect about 1); Cios Fusion only: split by view, no pooled interval.
\par\endgroup
\end{minipage}}
\endgroup
\end{table}
, so its table numbers S3). The supplement has no word limit (12 pages).
% Every value carries its source in the comment beside it;
% \PENDING marks what no file on disk yet states. Code SHA is absent from the level-model result files
% (docs/MODEL_CARD.md 1.2), so the training settings are those of the scripts named, at the revision on disk on
% 2026-09-23; paper 2's manifest restates them for the split-trained arm (the comparison is paper 2 since 2026-09-24,
% PAPER1_SCOPE.md).

\section*{Technical details}

\paragraph{Pre-processing.}
% docs/REAL_XRAY_DATASETS_SPEC.md "Display" and "Boxes"
Every real exposure was converted to attenuation (the negative logarithm of a linear detector signal against its
air or flood value, or a linear reading of a stored logarithmic display), then windowed identically for every
device: the field's 1st, 5th, 50th, 95th and 99th percentiles were mapped to grey levels 27, 47, 126, 216 and 239,
bone bright, and stored as an 8-bit image. A truth box is the projection of every voxel of a cone-beam CT vertebra
label (subsampled to at most 20,000 points) through the registered pose, taking the 2nd--98th percentile extent;
a box was kept when at least half its projected points fell inside the image and both of its sides measured at
least 8 pixels. Level crops are fixed-scale two-channel windows of 320 pixels (image and target marker) with the
target body at 96 pixels.

\paragraph{Acquisition.}
% docs/RSNAI_TABLE1_DATASETS.md (pitch, frame size, stored value, tube voltage p5-p95); docs/MODEL_CARD.md 6.1
% (C-ArmLSpine source-image and source-object distance). Retrospective public releases: no acquisition protocol
% beyond what the headers carry; "not recorded" where the release gives none.
Table~\ref{tab:acquisition} lists the acquisition parameters each release carries. No release documents an exposure protocol beyond
its headers, and no parameter was controlled by this study.

\begin{table}[tbh!]
\centering
\caption{Acquisition parameters of the real exposures, as released.\label{tab:acquisition}}
\footnotesize
\begin{tabular}{>{\raggedright\arraybackslash}p{2.5cm}>{\raggedright\arraybackslash}p{3.6cm}>{\raggedright\arraybackslash}p{1.7cm}>{\raggedright\arraybackslash}p{3.3cm}>{\raggedright\arraybackslash}p{2.6cm}}
\toprule
Collection & Device & Detector pitch (mm) & Frame (pixels) & Tube voltage, kVp (5th--95th percentile) \\
\midrule
C-ArmLSpine & Surgio B51 flat-panel C-arm & 0.15 (a) & 2,048 $\times$ 2,048; exported at 1,024 & 73--111 \\
DeepFluoro & Siemens Cios Fusion & 0.194 & 1,536 $\times$ 1,536 & not recorded (b) \\
COBRA2026 & Elekta XVI kV imager & 0.8 & 504 $\times$ 504 & 120 \\
Geneva hip & Philips BV Pulsera image intensifier & 0.2904 & 1,000 $\times$ 1,000 & 65--90 \\
Ljubljana 2D-3D-GS & Trixell PIXIUM 4600 flat panel & 0.286 & 1,560 $\times$ 1,560 & not recorded \\
\bottomrule
\end{tabular}

\smallskip
\raggedright\footnotesize
(a) 307.2 mm over 2,048 pixels; the header pixel spacing is a placeholder of 1.0. Source-to-image distance
1,140 mm, source-to-object distance 745 mm, object-plane pixel 0.1915 mm at the 1,024-pixel export.
(b) Present in the vendor DICOM headers of the same acquisitions but not joined per registered frame.
\end{table}

\paragraph{Models, initialisation and training.}
% Held-out-device detector: atlas/model.json "Architecture" (args.yaml and data.yaml of rx_lodo_carmlspine_1024 read
% on the grid 2026-09-22); results/a1.json operating_points/cross_device_0.74 (how_chosen).
The sacrum detector is YOLO11m (Ultralytics; classes vertebra and sacrum; 1,024-pixel input), initialised from the
COCO-pretrained YOLO11m weights through a render-trained checkpoint and trained for 80 epochs (patience 25, batch
8, seed 0, deterministic) with rotation up to 10\textdegree, translation 0.1, scale 0.3, horizontal flip 0.5 and
mosaic 0.5, on the training and validation frames of COBRA2026 and DeepFluoro; its checkpoint was selected on
their test frames, so only its held-out-device results are clean.
% Keypoint detector: docs/MANIFEST_DATA_MODELS_RESULTS.md row detector_v5z_final; Methods (confidence 0.35);
% chain detect_v5z_final <- detect_v3 <- public ultralytics/yolo11m-pose.pt (COCO-pose), results/eval_gate/LINEAGE.md.
The body detector is YOLO11m-pose fine-tuned for 60 epochs on the plain-film keypoint corpus from an earlier pose
checkpoint (detect\_v3), itself trained from the COCO-pretrained YOLO11m-pose weights, and used at confidence 0.35
\PENDING{G5}{keypoint detector: training corpus by source and split, and whether any film of a test person is in it}.
% Encoder: docs/MODEL_CARD.md 1.1 and 6.3 (architecture, 11,173,376 parameters, 3.62 GMAC); slurm/carmlspine_cv.sh
% lines 21 and 31 (init alias allspine_v2_r1, 15 epochs, lr 1e-4, --augment 2); tools/train_level_classifier.py
% (AdamW, weight decay 1e-4, cosine schedule, class-weighted cross-entropy).
The level encoder is ResNet-18 with a two-channel stem and a 27-level head (11,173,376 parameters; 3.62 giga
multiply--accumulates per crop). The cross-validation encoders were initialised from a level encoder trained on
renders and films (registry model allspine\_v2\_r1) and fine-tuned per fold for 15 epochs at a learning rate of
$10^{-4}$ with AdamW (weight decay $10^{-4}$), a cosine schedule and class-weighted cross-entropy.
% Figure 2 encoder: results/figure2/figure2_metrics.json "encoder", "encoder_sha256"; ImageNet root per
% results/eval_gate/LINEAGE.md (allspine_v9 <- public torchvision resnet18 IMAGENET1K_V1, "used when cached on the
% node; random init otherwise"). Checked 2026-09-24 (grid jobs 40429822, 40429828, read-only): the training script
% loads the ImageNet file when it is in the torch hub cache and prints "pretrained=<bool>"; the file was in the shared
% home cache from 2026-09-21 18:31, and v9 trained from 2026-09-22 21:57 (model_night/logs/v9train_40415669.out), but
% that log starts at epoch 25, so the load itself is not recorded. The sentence says so.
% Checkpoint selection: tools/train_level_classifier.py keeps the best epoch by exact accuracy on the bank's validation
% split (results/eval_gate/EVALUATORS.md item 1), which is why Table 3 (tab3_naming.tex; Table 4 before 2026-09-24
% evening) labels the film rows "selection split".
The encoder whose embedding Figure~2 maps is the allspine\_v9 model (SHA-256 188f9fcb\ldots), started by the training
script's default from ImageNet weights, which were cached on the computing grid before it trained (its log does not
record the load); like every level encoder here, its checkpoint was the epoch with the best exact accuracy on its
bank's validation split, so numbers on that split are selection-split numbers.
% Head: docs/MODEL_CARD.md 1.1 (RunLabeller, 2,117,403 parameters); tools/head_mixed_length.py (defaults epochs 30,
% lr 1e-3, AdamW weight decay 1e-4, cosine; per-position cross-entropy); slurm/carmlspine_cv.sh line 33
% (--train-mode mixed, no epoch or learning-rate override).
The recurrent head (linear 512 to 256, GELU, two-layer bidirectional GRU of width 256, linear 512 to 27; 2,117,403
parameters) was trained on runs of two to nine consecutive embeddings of mixed length for 30 epochs at a learning
rate of $10^{-3}$ (AdamW, weight decay $10^{-4}$, cosine schedule) with per-position cross-entropy.
% docs/MODEL_CARD.md 6.3 ("flip TTA on"); tools/alphabet_couple.py header (posterior "flip-TTA averaged"); no ensemble.
No ensemble is used; each crop is also read horizontally flipped and the two readings averaged.
Counting and the alphabet coupling have no learned parameters.
% Split-trained arm: docs/RESUME_2026-09-22_PM.md checkpoint 2026-09-23 21:45 (the v8/v9 banks trained on 49 of 54 L6
% validation and test persons of the built split, tables/tab4_transitional.md note d, so the arm starts from
% ImageNet on the split's training side). Since 2026-09-24 (evening) that arm is scored in paper 2 (the former
% "Figure 6 arm"; paper2_levelcheck/); paper 1 shows only its development learning curve (supplement section
% "Development learning curve of a split-trained level model"), so the sentence names no scored configuration.
A level model trained from ImageNet weights on the training side of the person-disjoint transitional split,
because the banks behind the encoders above hold most of its validation and test persons, is reported here only
through its development learning curve; no model in this manuscript is scored on a test person of that split.
 after
% % SUPPLEMENT Table S2 since the 2026-09-23 six-figure wrap (formerly main Table 2); input by supplement.tex.
% Table 2 -- detection per held-out device and training arm at both operating points, with case-clustered intervals
% and PR-AUC. No merged cells: the target block heading sits in the first cell of its first row. Sources, one line per
% cell, in tab2_detection.sources.md; the same rows with BCa intervals, plus the arms cut from the page for space
% (the design-A seed pairs, renders-only on COBRA and DeepFluoro, Cios-only on COBRA, Elekta-only and Surgio-only on
% DeepFluoro), in tab2_detection.md.
%   results/a1.json                  the reported detector (lodo_carm) on C-ArmLSpine at 0.74 and 0.25 (k/n, percentile, BCa)
%   results/lcurve/lcurve.json       design-B stage-1 arms and the same-device controls at 0.74 and 0.25, PR-AUC (both devices)
%   results/contribution_matrix.json every other arm at 0.25 (recall as a fraction; k = recall x n, exact integers)
%   results/a3.json                  the plain-film block (render-trained vs film-trained, conf 0.05 and 0.25)
%   results/fig3_detection/scored/*  the 0.74 cells of the render-only and real-only held-out arms (same sweeps, both points)
% BOX CONVENTION (A28): every number here is in the whole-vertebra convention, and the sentence in the text that quotes a
% row must say "whole-vertebra". The vertebral-body-box row per block (the all-device detector det_g3_primary_s0, 15 G5
% cells) was DROPPED from paper 1 on 2026-09-24 (Greg): it can be scored only on the held-out test persons paper 2 keeps
% blinded; the rows are kept verbatim in paper2_levelcheck/MOVED_FROM_PAPER1.md (item 10).
% Labels: tab:detection is the table; tab:contribution is kept as an alias so the Results prose that still cites the
% retired contribution table resolves to this one (the single-source arms it quoted are rows here).
% Revision 2 (2026-09-24, REVISION_LOG2.md): plain-radiograph rows re-scored on the 17 NHANES II lumbar spines neither
% detector trained or was selected on (EVAL_LEAKS L2; closure MANUSCRIPT_CORRECTIONS.md 3a;
% results/eval_gate/leaks/render_vs_film_unseen_rerun.json); zero-floor percentile cells given unclustered Wilson
% intervals (REVIEW2_STATS N18; results/stats_review2/recompute.json#S10_floor_cells).
% Layout: landscape (rotated 23.2 cm minipage, see tab_corpus.tex); body 7 pt on 7.6 pt.
\begin{table}[p]
\centering
\begingroup
\singlespacing
\fontsize{7}{7.6}\selectfont
\newcommand{\fnm}[1]{\textsuperscript{\fontsize{7}{7}\selectfont #1}}
\setlength{\tabcolsep}{2pt}
\renewcommand{\arraystretch}{1.0}
\rotatebox[origin=c]{90}{\begin{minipage}{23.2cm}
\captionsetup{font=footnotesize,skip=4pt}
\centering
\caption{Detection per held-out device and training arm at the cross-device (0.74) and default (0.25) operating points. Whole-vertebra box convention throughout; the all-device detector trained on vertebral-body boxes is reported in the separate held-out study.}
\label{tab:detection}\label{tab:contribution}
% widths sum to 22.00 cm; 14 gaps of 2 pt (0.99 cm) -> 22.99 cm
\begin{tabular}{%
>{\raggedright\arraybackslash}p{2.20cm}%
>{\raggedright\arraybackslash}p{5.20cm}%
>{\raggedright\arraybackslash}p{3.15cm}%
>{\raggedright\arraybackslash}p{3.00cm}%
>{\raggedright\arraybackslash}p{3.15cm}%
>{\raggedright\arraybackslash}p{3.00cm}%
>{\raggedright\arraybackslash}p{2.30cm}}
\toprule
Target device\fnm{a} & Training arm & Vertebra recall at 0.74\fnm{b} & Sacrum found at 0.74 & Vertebra recall at 0.25 & Sacrum found at 0.25 & PR-AUC\fnm{c}, vertebra / sacrum \\
\midrule
C-ArmLSpine (Surgio-B51) & Reported detector: real COBRA2026 and DeepFluoro frames, no render & 80.5 (231/287) [74.1, 86.8] & 76.8 (43/56) [64.3, 89.3] & 90.6 (260/287) [84.3, 94.5]\fnm{g} & 83.9 (47/56) [73.2, 92.9] & -- \\
 & Design B: styled renders and real frames of the two other devices; seed 1 & 63.1 (181/287) [54.8, 71.0] & 66.1 (37/56) [53.6, 78.6] & 87.8 (252/287) [81.4, 93.2] & 85.7 (48/56) [76.8, 92.9] & 0.890 / 0.900 \\
 & Elekta XVI only, styled and real & -- & -- & 86.8 (249/287) [81.1, 91.8] & 78.6 (44/56) [67.9, 89.3] & -- \\
 & Cios Fusion only, styled and real & -- & -- & 12.9 (37/287) [8.3, 18.1] & AP 28/28, lateral 0/28\fnm{g} & -- \\
 & Elekta XVI only, real frames cut to 240 in 10 patients & -- & -- & 39.0 (112/287) [33.2, 44.7] & 37.5 (21/56) [28.6, 46.4] & -- \\
 & Renders only, no real frame & 0.3 (1/287) [0.1, 1.9]\fnm{g} & 5.4 (3/56) [1.8, 14.6]\fnm{g} & 31.7 (91/287) [27.1, 36.1] & 53.6 (30/56) [44.6, 62.5] & -- \\
 & Control\fnm{d}: C-ArmLSpine real training frames only & 40.4 (116/287) [37.8, 42.9] & 94.6 (53/56) [85.4, 98.2]\fnm{g} & 41.5 (119/287) [38.8, 44.0] & 98.2 (55/56) [90.6, 99.7]\fnm{g} & 0.367 / 0.979 \\
 & Control\fnm{d}: real training frames of all three devices & 40.1 (115/287) [37.7, 42.3] & 96.4 (54/56) [87.9, 99.0]\fnm{g} & 40.8 (117/287) [38.4, 43.0] & 98.2 (55/56) [90.6, 99.7]\fnm{g} & 0.385 / 0.982 \\
\addlinespace[1.5pt]
COBRA2026 (Elekta XVI) & Real DeepFluoro and C-ArmLSpine frames, no render & 6.2 (55/893) [2.9, 9.4] & 3.4 (21/609) [1.3, 5.9] & 18.5 (165/893) [13.7, 22.8] & 13.3 (81/609) [8.1, 18.8] & -- \\
 & Design B: styled renders and real frames of the two other devices & 10.1 (90/893) [6.7, 13.1] & 13.0 (79/609) [9.0, 17.1] & 32.7 (292/893) [28.4, 36.2] & 36.5 (222/609) [31.6, 40.9] & 0.266 / 0.391 \\
 & Surgio-B51 only, styled and real; two runs of one list & -- & -- & 7.5 (67/893) [3.5, 12.1]; 12.2 (109/893) [8.0, 16.0] & 3.8 (23/609) [1.8, 6.2]; 8.0 (49/609) [5.3, 10.9] & -- \\
 & Control\fnm{d}: real training frames of all three devices & 93.2 (832/893) [86.4, 97.7] & 90.6 (552/609) [83.8, 96.3] & 97.4 (870/893) [93.3, 99.8] & 97.9 (596/609) [95.7, 99.5] & 0.963 / 0.985 \\
\addlinespace[1.5pt]
DeepFluoro (Cios Fusion)\fnm{e} & Real COBRA2026 and C-ArmLSpine frames, no render; test specimen & 74.7 (62/83)\fnm{e} & 100.0 (90/90)\fnm{e} & 95.2 (79/83)\fnm{e} & 100.0 (90/90)\fnm{e} & -- \\
 & Design B: styled renders and real frames of the two other devices; six specimens & -- & -- & 78.6 (338/430) [68.8, 92.0] & 95.0 (321/338) [78.5, 99.8] & -- \\
 & The same arm; test specimen & 69.9 (58/83)\fnm{e} & 87.8 (79/90)\fnm{e} & 94.0 (78/83)\fnm{e} & 100.0 (90/90)\fnm{e} & 0.923 / 0.954 \\
 & Control\fnm{d}: real training frames of all three devices; test specimen & 73.5 (61/83)\fnm{e} & 100.0 (90/90)\fnm{e} & 95.2 (79/83)\fnm{e} & 100.0 (90/90)\fnm{e} & 0.685 / 0.968 \\
\addlinespace[1.5pt]
Plain radiographs\fnm{f} & Render-trained detector, no real image & at 0.05: 67.7 (109/161) & -- & 34.8 (56/161) & -- & -- \\
 & Film-trained detector & at 0.05: 99.4 (160/161) & -- & 98.8 (159/161) & -- & -- \\
\bottomrule
\end{tabular}
\par\vspace{1pt}
\begingroup\fontsize{7}{7.6}\selectfont\raggedright
Cells are \% ($k/n$) [95\,\% CI]. CI = confidence interval; PR-AUC = area under the precision--recall curve; IoU = intersection over union. Greedy one-to-one matching at IoU $\ge$ 0.5; 95\,\% percentile bootstraps of 2,000 resamples clustered by case (C-ArmLSpine), patient (COBRA2026) or specimen (DeepFluoro), BCa intervals in the Markdown twin; ``--'' = not scored at that point or not computed. Styled = rendered through the fitted terms; design B = renders under each device's own poses.
\fnm{a}Test sets: C-ArmLSpine 56 frames, 28 cases, 287 vertebrae, 56 sacra; COBRA2026 648 frames, 27 patients, 893 vertebrae, 609 frames with a sacrum; DeepFluoro six specimens, 366 frames, 430 vertebrae, 338 with a sacrum; plain radiographs 17 NHANES II lumbar films (17 spines, 161 bodies) whose patients neither detector trained or was selected on. Truth boxes: C-ArmLSpine rows against whole-vertebra outline boxes projected from the CBCT labels, the convention every number on disk uses (under the SPINEPS vertebral-body boxes of the pending flagship rows the reported detector reads 5.2\,\% (15/287) at 0.25, a convention mismatch); COBRA2026 and DeepFluoro rows against whole-vertebra boxes projected from the registered CT. The device is held out unless the row says control.
\fnm{b}0.74 = best pooled vertebra F1 of the reported detector's sweeps on COBRA2026, DeepFluoro and the Geneva phantom, reading no C-ArmLSpine frame, adopted after the 0.25 result was reported; 0.25 = the software default; the reported detector chose its checkpoint on the COBRA2026 and DeepFluoro test splits.
\fnm{c}Non-interpolated average precision over the confidence sweep, computed for the design-B arms and the controls only.
\fnm{d}Controls trained on the target device's own frames (test cases held out by case); on C-ArmLSpine the outline convention penalises these body-trained arms, which read 89.9\,\% (258/287) and 90.6\,\% (260/287) at 0.25 under body boxes.
\fnm{e}One specimen is one cluster, no interval; the six-specimen rows test arms that never saw the device, their six-cluster intervals descriptive.
\fnm{f}Scored at 0.05 and 0.25 only, on every annotated NHANES~II lumbar spine that neither detector trained or was selected on; paired difference at 0.05 +31.7 percentage points [21.0, 45.4], sign test 16 vs 0 spines, Holm-adjusted $P<.001$; 17 clusters, so the intervals are approximate.
\fnm{g}Wilson 95\,\% interval where the percentile one reached 0 or 100\,\% against the data (unclustered for the two 0 cells) or the file flags errors in few clusters (sacrum controls unadjusted, design effect about 1); Cios Fusion only: split by view, no pooled interval.
\par\endgroup
\end{minipage}}
\endgroup
\end{table}
, so its table numbers S3). The supplement has no word limit (12 pages).
% Every value carries its source in the comment beside it;
% \PENDING marks what no file on disk yet states. Code SHA is absent from the level-model result files
% (docs/MODEL_CARD.md 1.2), so the training settings are those of the scripts named, at the revision on disk on
% 2026-09-23; paper 2's manifest restates them for the split-trained arm (the comparison is paper 2 since 2026-09-24,
% PAPER1_SCOPE.md).

\section*{Technical details}

\paragraph{Pre-processing.}
% docs/REAL_XRAY_DATASETS_SPEC.md "Display" and "Boxes"
Every real exposure was converted to attenuation (the negative logarithm of a linear detector signal against its
air or flood value, or a linear reading of a stored logarithmic display), then windowed identically for every
device: the field's 1st, 5th, 50th, 95th and 99th percentiles were mapped to grey levels 27, 47, 126, 216 and 239,
bone bright, and stored as an 8-bit image. A truth box is the projection of every voxel of a cone-beam CT vertebra
label (subsampled to at most 20,000 points) through the registered pose, taking the 2nd--98th percentile extent;
a box was kept when at least half its projected points fell inside the image and both of its sides measured at
least 8 pixels. Level crops are fixed-scale two-channel windows of 320 pixels (image and target marker) with the
target body at 96 pixels.

\paragraph{Acquisition.}
% docs/RSNAI_TABLE1_DATASETS.md (pitch, frame size, stored value, tube voltage p5-p95); docs/MODEL_CARD.md 6.1
% (C-ArmLSpine source-image and source-object distance). Retrospective public releases: no acquisition protocol
% beyond what the headers carry; "not recorded" where the release gives none.
Table~\ref{tab:acquisition} lists the acquisition parameters each release carries. No release documents an exposure protocol beyond
its headers, and no parameter was controlled by this study.

\begin{table}[tbh!]
\centering
\caption{Acquisition parameters of the real exposures, as released.\label{tab:acquisition}}
\footnotesize
\begin{tabular}{>{\raggedright\arraybackslash}p{2.5cm}>{\raggedright\arraybackslash}p{3.6cm}>{\raggedright\arraybackslash}p{1.7cm}>{\raggedright\arraybackslash}p{3.3cm}>{\raggedright\arraybackslash}p{2.6cm}}
\toprule
Collection & Device & Detector pitch (mm) & Frame (pixels) & Tube voltage, kVp (5th--95th percentile) \\
\midrule
C-ArmLSpine & Surgio B51 flat-panel C-arm & 0.15 (a) & 2,048 $\times$ 2,048; exported at 1,024 & 73--111 \\
DeepFluoro & Siemens Cios Fusion & 0.194 & 1,536 $\times$ 1,536 & not recorded (b) \\
COBRA2026 & Elekta XVI kV imager & 0.8 & 504 $\times$ 504 & 120 \\
Geneva hip & Philips BV Pulsera image intensifier & 0.2904 & 1,000 $\times$ 1,000 & 65--90 \\
Ljubljana 2D-3D-GS & Trixell PIXIUM 4600 flat panel & 0.286 & 1,560 $\times$ 1,560 & not recorded \\
\bottomrule
\end{tabular}

\smallskip
\raggedright\footnotesize
(a) 307.2 mm over 2,048 pixels; the header pixel spacing is a placeholder of 1.0. Source-to-image distance
1,140 mm, source-to-object distance 745 mm, object-plane pixel 0.1915 mm at the 1,024-pixel export.
(b) Present in the vendor DICOM headers of the same acquisitions but not joined per registered frame.
\end{table}

\paragraph{Models, initialisation and training.}
% Held-out-device detector: atlas/model.json "Architecture" (args.yaml and data.yaml of rx_lodo_carmlspine_1024 read
% on the grid 2026-09-22); results/a1.json operating_points/cross_device_0.74 (how_chosen).
The sacrum detector is YOLO11m (Ultralytics; classes vertebra and sacrum; 1,024-pixel input), initialised from the
COCO-pretrained YOLO11m weights through a render-trained checkpoint and trained for 80 epochs (patience 25, batch
8, seed 0, deterministic) with rotation up to 10\textdegree, translation 0.1, scale 0.3, horizontal flip 0.5 and
mosaic 0.5, on the training and validation frames of COBRA2026 and DeepFluoro; its checkpoint was selected on
their test frames, so only its held-out-device results are clean.
% Keypoint detector: docs/MANIFEST_DATA_MODELS_RESULTS.md row detector_v5z_final; Methods (confidence 0.35);
% chain detect_v5z_final <- detect_v3 <- public ultralytics/yolo11m-pose.pt (COCO-pose), results/eval_gate/LINEAGE.md.
The body detector is YOLO11m-pose fine-tuned for 60 epochs on the plain-film keypoint corpus from an earlier pose
checkpoint (detect\_v3), itself trained from the COCO-pretrained YOLO11m-pose weights, and used at confidence 0.35
\PENDING{G5}{keypoint detector: training corpus by source and split, and whether any film of a test person is in it}.
% Encoder: docs/MODEL_CARD.md 1.1 and 6.3 (architecture, 11,173,376 parameters, 3.62 GMAC); slurm/carmlspine_cv.sh
% lines 21 and 31 (init alias allspine_v2_r1, 15 epochs, lr 1e-4, --augment 2); tools/train_level_classifier.py
% (AdamW, weight decay 1e-4, cosine schedule, class-weighted cross-entropy).
The level encoder is ResNet-18 with a two-channel stem and a 27-level head (11,173,376 parameters; 3.62 giga
multiply--accumulates per crop). The cross-validation encoders were initialised from a level encoder trained on
renders and films (registry model allspine\_v2\_r1) and fine-tuned per fold for 15 epochs at a learning rate of
$10^{-4}$ with AdamW (weight decay $10^{-4}$), a cosine schedule and class-weighted cross-entropy.
% Figure 2 encoder: results/figure2/figure2_metrics.json "encoder", "encoder_sha256"; ImageNet root per
% results/eval_gate/LINEAGE.md (allspine_v9 <- public torchvision resnet18 IMAGENET1K_V1, "used when cached on the
% node; random init otherwise"). Checked 2026-09-24 (grid jobs 40429822, 40429828, read-only): the training script
% loads the ImageNet file when it is in the torch hub cache and prints "pretrained=<bool>"; the file was in the shared
% home cache from 2026-09-21 18:31, and v9 trained from 2026-09-22 21:57 (model_night/logs/v9train_40415669.out), but
% that log starts at epoch 25, so the load itself is not recorded. The sentence says so.
% Checkpoint selection: tools/train_level_classifier.py keeps the best epoch by exact accuracy on the bank's validation
% split (results/eval_gate/EVALUATORS.md item 1), which is why Table 3 (tab3_naming.tex; Table 4 before 2026-09-24
% evening) labels the film rows "selection split".
The encoder whose embedding Figure~2 maps is the allspine\_v9 model (SHA-256 188f9fcb\ldots), started by the training
script's default from ImageNet weights, which were cached on the computing grid before it trained (its log does not
record the load); like every level encoder here, its checkpoint was the epoch with the best exact accuracy on its
bank's validation split, so numbers on that split are selection-split numbers.
% Head: docs/MODEL_CARD.md 1.1 (RunLabeller, 2,117,403 parameters); tools/head_mixed_length.py (defaults epochs 30,
% lr 1e-3, AdamW weight decay 1e-4, cosine; per-position cross-entropy); slurm/carmlspine_cv.sh line 33
% (--train-mode mixed, no epoch or learning-rate override).
The recurrent head (linear 512 to 256, GELU, two-layer bidirectional GRU of width 256, linear 512 to 27; 2,117,403
parameters) was trained on runs of two to nine consecutive embeddings of mixed length for 30 epochs at a learning
rate of $10^{-3}$ (AdamW, weight decay $10^{-4}$, cosine schedule) with per-position cross-entropy.
% docs/MODEL_CARD.md 6.3 ("flip TTA on"); tools/alphabet_couple.py header (posterior "flip-TTA averaged"); no ensemble.
No ensemble is used; each crop is also read horizontally flipped and the two readings averaged.
Counting and the alphabet coupling have no learned parameters.
% Split-trained arm: docs/RESUME_2026-09-22_PM.md checkpoint 2026-09-23 21:45 (the v8/v9 banks trained on 49 of 54 L6
% validation and test persons of the built split, tables/tab4_transitional.md note d, so the arm starts from
% ImageNet on the split's training side). Since 2026-09-24 (evening) that arm is scored in paper 2 (the former
% "Figure 6 arm"; paper2_levelcheck/); paper 1 shows only its development learning curve (supplement section
% "Development learning curve of a split-trained level model"), so the sentence names no scored configuration.
A level model trained from ImageNet weights on the training side of the person-disjoint transitional split,
because the banks behind the encoders above hold most of its validation and test persons, is reported here only
through its development learning curve; no model in this manuscript is scored on a test person of that split.
 after
% % SUPPLEMENT Table S2 since the 2026-09-23 six-figure wrap (formerly main Table 2); input by supplement.tex.
% Table 2 -- detection per held-out device and training arm at both operating points, with case-clustered intervals
% and PR-AUC. No merged cells: the target block heading sits in the first cell of its first row. Sources, one line per
% cell, in tab2_detection.sources.md; the same rows with BCa intervals, plus the arms cut from the page for space
% (the design-A seed pairs, renders-only on COBRA and DeepFluoro, Cios-only on COBRA, Elekta-only and Surgio-only on
% DeepFluoro), in tab2_detection.md.
%   results/a1.json                  the reported detector (lodo_carm) on C-ArmLSpine at 0.74 and 0.25 (k/n, percentile, BCa)
%   results/lcurve/lcurve.json       design-B stage-1 arms and the same-device controls at 0.74 and 0.25, PR-AUC (both devices)
%   results/contribution_matrix.json every other arm at 0.25 (recall as a fraction; k = recall x n, exact integers)
%   results/a3.json                  the plain-film block (render-trained vs film-trained, conf 0.05 and 0.25)
%   results/fig3_detection/scored/*  the 0.74 cells of the render-only and real-only held-out arms (same sweeps, both points)
% BOX CONVENTION (A28): every number here is in the whole-vertebra convention, and the sentence in the text that quotes a
% row must say "whole-vertebra". The vertebral-body-box row per block (the all-device detector det_g3_primary_s0, 15 G5
% cells) was DROPPED from paper 1 on 2026-09-24 (Greg): it can be scored only on the held-out test persons paper 2 keeps
% blinded; the rows are kept verbatim in paper2_levelcheck/MOVED_FROM_PAPER1.md (item 10).
% Labels: tab:detection is the table; tab:contribution is kept as an alias so the Results prose that still cites the
% retired contribution table resolves to this one (the single-source arms it quoted are rows here).
% Revision 2 (2026-09-24, REVISION_LOG2.md): plain-radiograph rows re-scored on the 17 NHANES II lumbar spines neither
% detector trained or was selected on (EVAL_LEAKS L2; closure MANUSCRIPT_CORRECTIONS.md 3a;
% results/eval_gate/leaks/render_vs_film_unseen_rerun.json); zero-floor percentile cells given unclustered Wilson
% intervals (REVIEW2_STATS N18; results/stats_review2/recompute.json#S10_floor_cells).
% Layout: landscape (rotated 23.2 cm minipage, see tab_corpus.tex); body 7 pt on 7.6 pt.
\begin{table}[p]
\centering
\begingroup
\singlespacing
\fontsize{7}{7.6}\selectfont
\newcommand{\fnm}[1]{\textsuperscript{\fontsize{7}{7}\selectfont #1}}
\setlength{\tabcolsep}{2pt}
\renewcommand{\arraystretch}{1.0}
\rotatebox[origin=c]{90}{\begin{minipage}{23.2cm}
\captionsetup{font=footnotesize,skip=4pt}
\centering
\caption{Detection per held-out device and training arm at the cross-device (0.74) and default (0.25) operating points. Whole-vertebra box convention throughout; the all-device detector trained on vertebral-body boxes is reported in the separate held-out study.}
\label{tab:detection}\label{tab:contribution}
% widths sum to 22.00 cm; 14 gaps of 2 pt (0.99 cm) -> 22.99 cm
\begin{tabular}{%
>{\raggedright\arraybackslash}p{2.20cm}%
>{\raggedright\arraybackslash}p{5.20cm}%
>{\raggedright\arraybackslash}p{3.15cm}%
>{\raggedright\arraybackslash}p{3.00cm}%
>{\raggedright\arraybackslash}p{3.15cm}%
>{\raggedright\arraybackslash}p{3.00cm}%
>{\raggedright\arraybackslash}p{2.30cm}}
\toprule
Target device\fnm{a} & Training arm & Vertebra recall at 0.74\fnm{b} & Sacrum found at 0.74 & Vertebra recall at 0.25 & Sacrum found at 0.25 & PR-AUC\fnm{c}, vertebra / sacrum \\
\midrule
C-ArmLSpine (Surgio-B51) & Reported detector: real COBRA2026 and DeepFluoro frames, no render & 80.5 (231/287) [74.1, 86.8] & 76.8 (43/56) [64.3, 89.3] & 90.6 (260/287) [84.3, 94.5]\fnm{g} & 83.9 (47/56) [73.2, 92.9] & -- \\
 & Design B: styled renders and real frames of the two other devices; seed 1 & 63.1 (181/287) [54.8, 71.0] & 66.1 (37/56) [53.6, 78.6] & 87.8 (252/287) [81.4, 93.2] & 85.7 (48/56) [76.8, 92.9] & 0.890 / 0.900 \\
 & Elekta XVI only, styled and real & -- & -- & 86.8 (249/287) [81.1, 91.8] & 78.6 (44/56) [67.9, 89.3] & -- \\
 & Cios Fusion only, styled and real & -- & -- & 12.9 (37/287) [8.3, 18.1] & AP 28/28, lateral 0/28\fnm{g} & -- \\
 & Elekta XVI only, real frames cut to 240 in 10 patients & -- & -- & 39.0 (112/287) [33.2, 44.7] & 37.5 (21/56) [28.6, 46.4] & -- \\
 & Renders only, no real frame & 0.3 (1/287) [0.1, 1.9]\fnm{g} & 5.4 (3/56) [1.8, 14.6]\fnm{g} & 31.7 (91/287) [27.1, 36.1] & 53.6 (30/56) [44.6, 62.5] & -- \\
 & Control\fnm{d}: C-ArmLSpine real training frames only & 40.4 (116/287) [37.8, 42.9] & 94.6 (53/56) [85.4, 98.2]\fnm{g} & 41.5 (119/287) [38.8, 44.0] & 98.2 (55/56) [90.6, 99.7]\fnm{g} & 0.367 / 0.979 \\
 & Control\fnm{d}: real training frames of all three devices & 40.1 (115/287) [37.7, 42.3] & 96.4 (54/56) [87.9, 99.0]\fnm{g} & 40.8 (117/287) [38.4, 43.0] & 98.2 (55/56) [90.6, 99.7]\fnm{g} & 0.385 / 0.982 \\
\addlinespace[1.5pt]
COBRA2026 (Elekta XVI) & Real DeepFluoro and C-ArmLSpine frames, no render & 6.2 (55/893) [2.9, 9.4] & 3.4 (21/609) [1.3, 5.9] & 18.5 (165/893) [13.7, 22.8] & 13.3 (81/609) [8.1, 18.8] & -- \\
 & Design B: styled renders and real frames of the two other devices & 10.1 (90/893) [6.7, 13.1] & 13.0 (79/609) [9.0, 17.1] & 32.7 (292/893) [28.4, 36.2] & 36.5 (222/609) [31.6, 40.9] & 0.266 / 0.391 \\
 & Surgio-B51 only, styled and real; two runs of one list & -- & -- & 7.5 (67/893) [3.5, 12.1]; 12.2 (109/893) [8.0, 16.0] & 3.8 (23/609) [1.8, 6.2]; 8.0 (49/609) [5.3, 10.9] & -- \\
 & Control\fnm{d}: real training frames of all three devices & 93.2 (832/893) [86.4, 97.7] & 90.6 (552/609) [83.8, 96.3] & 97.4 (870/893) [93.3, 99.8] & 97.9 (596/609) [95.7, 99.5] & 0.963 / 0.985 \\
\addlinespace[1.5pt]
DeepFluoro (Cios Fusion)\fnm{e} & Real COBRA2026 and C-ArmLSpine frames, no render; test specimen & 74.7 (62/83)\fnm{e} & 100.0 (90/90)\fnm{e} & 95.2 (79/83)\fnm{e} & 100.0 (90/90)\fnm{e} & -- \\
 & Design B: styled renders and real frames of the two other devices; six specimens & -- & -- & 78.6 (338/430) [68.8, 92.0] & 95.0 (321/338) [78.5, 99.8] & -- \\
 & The same arm; test specimen & 69.9 (58/83)\fnm{e} & 87.8 (79/90)\fnm{e} & 94.0 (78/83)\fnm{e} & 100.0 (90/90)\fnm{e} & 0.923 / 0.954 \\
 & Control\fnm{d}: real training frames of all three devices; test specimen & 73.5 (61/83)\fnm{e} & 100.0 (90/90)\fnm{e} & 95.2 (79/83)\fnm{e} & 100.0 (90/90)\fnm{e} & 0.685 / 0.968 \\
\addlinespace[1.5pt]
Plain radiographs\fnm{f} & Render-trained detector, no real image & at 0.05: 67.7 (109/161) & -- & 34.8 (56/161) & -- & -- \\
 & Film-trained detector & at 0.05: 99.4 (160/161) & -- & 98.8 (159/161) & -- & -- \\
\bottomrule
\end{tabular}
\par\vspace{1pt}
\begingroup\fontsize{7}{7.6}\selectfont\raggedright
Cells are \% ($k/n$) [95\,\% CI]. CI = confidence interval; PR-AUC = area under the precision--recall curve; IoU = intersection over union. Greedy one-to-one matching at IoU $\ge$ 0.5; 95\,\% percentile bootstraps of 2,000 resamples clustered by case (C-ArmLSpine), patient (COBRA2026) or specimen (DeepFluoro), BCa intervals in the Markdown twin; ``--'' = not scored at that point or not computed. Styled = rendered through the fitted terms; design B = renders under each device's own poses.
\fnm{a}Test sets: C-ArmLSpine 56 frames, 28 cases, 287 vertebrae, 56 sacra; COBRA2026 648 frames, 27 patients, 893 vertebrae, 609 frames with a sacrum; DeepFluoro six specimens, 366 frames, 430 vertebrae, 338 with a sacrum; plain radiographs 17 NHANES II lumbar films (17 spines, 161 bodies) whose patients neither detector trained or was selected on. Truth boxes: C-ArmLSpine rows against whole-vertebra outline boxes projected from the CBCT labels, the convention every number on disk uses (under the SPINEPS vertebral-body boxes of the pending flagship rows the reported detector reads 5.2\,\% (15/287) at 0.25, a convention mismatch); COBRA2026 and DeepFluoro rows against whole-vertebra boxes projected from the registered CT. The device is held out unless the row says control.
\fnm{b}0.74 = best pooled vertebra F1 of the reported detector's sweeps on COBRA2026, DeepFluoro and the Geneva phantom, reading no C-ArmLSpine frame, adopted after the 0.25 result was reported; 0.25 = the software default; the reported detector chose its checkpoint on the COBRA2026 and DeepFluoro test splits.
\fnm{c}Non-interpolated average precision over the confidence sweep, computed for the design-B arms and the controls only.
\fnm{d}Controls trained on the target device's own frames (test cases held out by case); on C-ArmLSpine the outline convention penalises these body-trained arms, which read 89.9\,\% (258/287) and 90.6\,\% (260/287) at 0.25 under body boxes.
\fnm{e}One specimen is one cluster, no interval; the six-specimen rows test arms that never saw the device, their six-cluster intervals descriptive.
\fnm{f}Scored at 0.05 and 0.25 only, on every annotated NHANES~II lumbar spine that neither detector trained or was selected on; paired difference at 0.05 +31.7 percentage points [21.0, 45.4], sign test 16 vs 0 spines, Holm-adjusted $P<.001$; 17 clusters, so the intervals are approximate.
\fnm{g}Wilson 95\,\% interval where the percentile one reached 0 or 100\,\% against the data (unclustered for the two 0 cells) or the file flags errors in few clusters (sacrum controls unadjusted, design effect about 1); Cios Fusion only: split by view, no pooled interval.
\par\endgroup
\end{minipage}}
\endgroup
\end{table}
, so its table numbers S3). The supplement has no word limit (12 pages).
% Every value carries its source in the comment beside it;
% \PENDING marks what no file on disk yet states. Code SHA is absent from the level-model result files
% (docs/MODEL_CARD.md 1.2), so the training settings are those of the scripts named, at the revision on disk on
% 2026-09-23; paper 2's manifest restates them for the split-trained arm (the comparison is paper 2 since 2026-09-24,
% PAPER1_SCOPE.md).

\section*{Technical details}

\paragraph{Pre-processing.}
% docs/REAL_XRAY_DATASETS_SPEC.md "Display" and "Boxes"
Every real exposure was converted to attenuation (the negative logarithm of a linear detector signal against its
air or flood value, or a linear reading of a stored logarithmic display), then windowed identically for every
device: the field's 1st, 5th, 50th, 95th and 99th percentiles were mapped to grey levels 27, 47, 126, 216 and 239,
bone bright, and stored as an 8-bit image. A truth box is the projection of every voxel of a cone-beam CT vertebra
label (subsampled to at most 20,000 points) through the registered pose, taking the 2nd--98th percentile extent;
a box was kept when at least half its projected points fell inside the image and both of its sides measured at
least 8 pixels. Level crops are fixed-scale two-channel windows of 320 pixels (image and target marker) with the
target body at 96 pixels.

\paragraph{Acquisition.}
% docs/RSNAI_TABLE1_DATASETS.md (pitch, frame size, stored value, tube voltage p5-p95); docs/MODEL_CARD.md 6.1
% (C-ArmLSpine source-image and source-object distance). Retrospective public releases: no acquisition protocol
% beyond what the headers carry; "not recorded" where the release gives none.
Table~\ref{tab:acquisition} lists the acquisition parameters each release carries. No release documents an exposure protocol beyond
its headers, and no parameter was controlled by this study.

\begin{table}[tbh!]
\centering
\caption{Acquisition parameters of the real exposures, as released.\label{tab:acquisition}}
\footnotesize
\begin{tabular}{>{\raggedright\arraybackslash}p{2.5cm}>{\raggedright\arraybackslash}p{3.6cm}>{\raggedright\arraybackslash}p{1.7cm}>{\raggedright\arraybackslash}p{3.3cm}>{\raggedright\arraybackslash}p{2.6cm}}
\toprule
Collection & Device & Detector pitch (mm) & Frame (pixels) & Tube voltage, kVp (5th--95th percentile) \\
\midrule
C-ArmLSpine & Surgio B51 flat-panel C-arm & 0.15 (a) & 2,048 $\times$ 2,048; exported at 1,024 & 73--111 \\
DeepFluoro & Siemens Cios Fusion & 0.194 & 1,536 $\times$ 1,536 & not recorded (b) \\
COBRA2026 & Elekta XVI kV imager & 0.8 & 504 $\times$ 504 & 120 \\
Geneva hip & Philips BV Pulsera image intensifier & 0.2904 & 1,000 $\times$ 1,000 & 65--90 \\
Ljubljana 2D-3D-GS & Trixell PIXIUM 4600 flat panel & 0.286 & 1,560 $\times$ 1,560 & not recorded \\
\bottomrule
\end{tabular}

\smallskip
\raggedright\footnotesize
(a) 307.2 mm over 2,048 pixels; the header pixel spacing is a placeholder of 1.0. Source-to-image distance
1,140 mm, source-to-object distance 745 mm, object-plane pixel 0.1915 mm at the 1,024-pixel export.
(b) Present in the vendor DICOM headers of the same acquisitions but not joined per registered frame.
\end{table}

\paragraph{Models, initialisation and training.}
% Held-out-device detector: atlas/model.json "Architecture" (args.yaml and data.yaml of rx_lodo_carmlspine_1024 read
% on the grid 2026-09-22); results/a1.json operating_points/cross_device_0.74 (how_chosen).
The sacrum detector is YOLO11m (Ultralytics; classes vertebra and sacrum; 1,024-pixel input), initialised from the
COCO-pretrained YOLO11m weights through a render-trained checkpoint and trained for 80 epochs (patience 25, batch
8, seed 0, deterministic) with rotation up to 10\textdegree, translation 0.1, scale 0.3, horizontal flip 0.5 and
mosaic 0.5, on the training and validation frames of COBRA2026 and DeepFluoro; its checkpoint was selected on
their test frames, so only its held-out-device results are clean.
% Keypoint detector: docs/MANIFEST_DATA_MODELS_RESULTS.md row detector_v5z_final; Methods (confidence 0.35);
% chain detect_v5z_final <- detect_v3 <- public ultralytics/yolo11m-pose.pt (COCO-pose), results/eval_gate/LINEAGE.md.
The body detector is YOLO11m-pose fine-tuned for 60 epochs on the plain-film keypoint corpus from an earlier pose
checkpoint (detect\_v3), itself trained from the COCO-pretrained YOLO11m-pose weights, and used at confidence 0.35
\PENDING{G5}{keypoint detector: training corpus by source and split, and whether any film of a test person is in it}.
% Encoder: docs/MODEL_CARD.md 1.1 and 6.3 (architecture, 11,173,376 parameters, 3.62 GMAC); slurm/carmlspine_cv.sh
% lines 21 and 31 (init alias allspine_v2_r1, 15 epochs, lr 1e-4, --augment 2); tools/train_level_classifier.py
% (AdamW, weight decay 1e-4, cosine schedule, class-weighted cross-entropy).
The level encoder is ResNet-18 with a two-channel stem and a 27-level head (11,173,376 parameters; 3.62 giga
multiply--accumulates per crop). The cross-validation encoders were initialised from a level encoder trained on
renders and films (registry model allspine\_v2\_r1) and fine-tuned per fold for 15 epochs at a learning rate of
$10^{-4}$ with AdamW (weight decay $10^{-4}$), a cosine schedule and class-weighted cross-entropy.
% Figure 2 encoder: results/figure2/figure2_metrics.json "encoder", "encoder_sha256"; ImageNet root per
% results/eval_gate/LINEAGE.md (allspine_v9 <- public torchvision resnet18 IMAGENET1K_V1, "used when cached on the
% node; random init otherwise"). Checked 2026-09-24 (grid jobs 40429822, 40429828, read-only): the training script
% loads the ImageNet file when it is in the torch hub cache and prints "pretrained=<bool>"; the file was in the shared
% home cache from 2026-09-21 18:31, and v9 trained from 2026-09-22 21:57 (model_night/logs/v9train_40415669.out), but
% that log starts at epoch 25, so the load itself is not recorded. The sentence says so.
% Checkpoint selection: tools/train_level_classifier.py keeps the best epoch by exact accuracy on the bank's validation
% split (results/eval_gate/EVALUATORS.md item 1), which is why Table 3 (tab3_naming.tex; Table 4 before 2026-09-24
% evening) labels the film rows "selection split".
The encoder whose embedding Figure~2 maps is the allspine\_v9 model (SHA-256 188f9fcb\ldots), started by the training
script's default from ImageNet weights, which were cached on the computing grid before it trained (its log does not
record the load); like every level encoder here, its checkpoint was the epoch with the best exact accuracy on its
bank's validation split, so numbers on that split are selection-split numbers.
% Head: docs/MODEL_CARD.md 1.1 (RunLabeller, 2,117,403 parameters); tools/head_mixed_length.py (defaults epochs 30,
% lr 1e-3, AdamW weight decay 1e-4, cosine; per-position cross-entropy); slurm/carmlspine_cv.sh line 33
% (--train-mode mixed, no epoch or learning-rate override).
The recurrent head (linear 512 to 256, GELU, two-layer bidirectional GRU of width 256, linear 512 to 27; 2,117,403
parameters) was trained on runs of two to nine consecutive embeddings of mixed length for 30 epochs at a learning
rate of $10^{-3}$ (AdamW, weight decay $10^{-4}$, cosine schedule) with per-position cross-entropy.
% docs/MODEL_CARD.md 6.3 ("flip TTA on"); tools/alphabet_couple.py header (posterior "flip-TTA averaged"); no ensemble.
No ensemble is used; each crop is also read horizontally flipped and the two readings averaged.
Counting and the alphabet coupling have no learned parameters.
% Split-trained arm: docs/RESUME_2026-09-22_PM.md checkpoint 2026-09-23 21:45 (the v8/v9 banks trained on 49 of 54 L6
% validation and test persons of the built split, tables/tab4_transitional.md note d, so the arm starts from
% ImageNet on the split's training side). Since 2026-09-24 (evening) that arm is scored in paper 2 (the former
% "Figure 6 arm"; paper2_levelcheck/); paper 1 shows only its development learning curve (supplement section
% "Development learning curve of a split-trained level model"), so the sentence names no scored configuration.
A level model trained from ImageNet weights on the training side of the person-disjoint transitional split,
because the banks behind the encoders above hold most of its validation and test persons, is reported here only
through its development learning curve; no model in this manuscript is scored on a test person of that split.
